\documentclass[11pt,a4paper]{article}
\usepackage[utf8]{inputenc}
\usepackage[italian]{babel}
\usepackage{geometry}
\usepackage{titlesec}
\usepackage{enumitem}
\usepackage{booktabs}
\usepackage{tabularx}
\usepackage{color}
\usepackage{colortbl}
\usepackage{tcolorbox}
\usepackage{fancyhdr}
\usepackage{graphicx}

\geometry{margin=1.5cm}

\definecolor{militarygray}{RGB}{50, 50, 60}
\definecolor{militarylightgray}{RGB}{80, 80, 100}
\definecolor{militaryred}{RGB}{150, 30, 30}

\titleformat{\section}
  {\normalfont\Large\bfseries\color{militaryred}}
  {}{0em}{}[\titlerule]

\titleformat{\subsection}
  {\normalfont\large\bfseries\color{militaryred}}
  {}{0em}{}

\pagestyle{fancy}
\fancyhf{}
\renewcommand{\headrulewidth}{1pt}
\renewcommand{\footrulewidth}{1pt}
\fancyhead[L]{HSS Ironclad}
\fancyhead[C]{Nave Militare}
\fancyhead[R]{Coalizione}
\fancyfoot[C]{\thepage}

\begin{document}

\begin{center}
\Huge\textbf{\textcolor{militaryred}{HSS Ironclad (Nave Militare)}}\\[0.5cm]
\Large\textbf{Pannello Master Diceless}
\end{center}

\begin{tcolorbox}[colback=militarylightgray!20, colframe=militaryred, title=\textbf{Informazioni sulla Nave}]
\begin{tabular}{ll}
\textbf{Tipo}: & Incrociatore pesante della Coalizione \\
\textbf{Equipaggio}: & 150 membri (personale militare diversificato) \\
\textbf{Armamento}: & Armamento pesante, sistemi difensivi avanzati \\
\textbf{Specializzazione}: & Operazioni di sicurezza, difesa tattica, neutralizzazione minacce \\
\textbf{Origine}: & Coalizione dei Mondi Indipendenti della Shackleton Expanse \\
\end{tabular}
\end{tcolorbox}

\section{Quantum Data Points (QDP)}

\begin{tcolorbox}[colback=militarylightgray!10, colframe=militaryred]
\begin{tabularx}{\textwidth}{|X|X|X|}
\hline
\rowcolor{militaryred!20} \textbf{QDP Alpha} & \textbf{QDP Beta} & \textbf{QDP Gamma} \\
\hline
\large $\square\square$ \textbf{0} & \large $\square\square$ \textbf{0} & \large $\square\square$ \textbf{2} \\
\hline
\textbf{Conoscenza, Analisi,} & \textbf{Tattiche, Diplomazia,} & \textbf{Sistemi Difensivi,} \\
\textbf{Scienza} & \textbf{Sistemi Offensivi} & \textbf{Materiali Resistenti} \\
\hline
\end{tabularx}
\end{tcolorbox}

\section{Filosofia di Gioco Diceless}

\begin{tcolorbox}[colback=militarylightgray!10, colframe=militaryred]
\begin{itemize}[leftmargin=*]
\item Risoluzione narrativa basata sulla logica e sulla coerenza della storia
\item Competenza dei personaggi come presupposto - sono professionisti militari
\item Le scelte e le strategie hanno più peso dei tiri casuali di dadi
\item Tensione creata da dilemmi morali e scelte difficili, non dal rischio di fallimento
\item Le competenze sono autorizzazioni ad agire in determinati ambiti
\end{itemize}
\end{tcolorbox}

\section{La Coalizione dei Mondi Indipendenti}

\begin{tcolorbox}[colback=militarylightgray!10, colframe=militaryred]
\begin{itemize}[leftmargin=*]
\item Alleanza di colonie e mondi che cercano indipendenza dalle grandi potenze
\item Cultura basata su pragmatismo, autosufficienza e difesa condivisa
\item Diffidenza verso la Federazione, i Klingon e i Romulani
\item Forte etica di autodifesa e preparazione al peggio
\item Tecnologia ibrida derivata da diverse fonti e adattata alle loro necessità
\item Valorizzazione della diversità culturale come punto di forza strategica
\end{itemize}
\end{tcolorbox}

\section{Fasi dell'Avventura}

\begin{tcolorbox}[colback=militarylightgray!10, colframe=militaryred]
\begin{enumerate}[leftmargin=*]
\item \textbf{Introduzione (15 minuti)}: Stabilire perimetro di sicurezza attorno all'anomalia
\item \textbf{Missioni Iniziali (60 minuti)}: "Perimetro di Sicurezza" - Monitoraggio dell'anomalia e delle altre navi
\item \textbf{Svolta (90 minuti)}: "Segnale di Targetizzazione" - L'artefatto sembra puntare verso colonie della Coalizione
\item \textbf{Conclusione (60 minuti)}: Decisione finale sulla neutralizzazione o utilizzo dell'artefatto
\item \textbf{Concilio Iconiano (15 minuti)}: Negoziazione finale con le altre fazioni per determinare l'esito collettivo
\end{enumerate}
\end{tcolorbox}

\section{Enigmi della HSS Ironclad}

\begin{tcolorbox}[colback=militarylightgray!10, colframe=militaryred]
\begin{description}[leftmargin=*]
\item[Il Codice del Guerriero] I sistemi d'arma pulsano con energia in pattern che corrispondono ad antiche tattiche di battaglia iconiane.
\item[Le Voci dei Mondi Perduti] L'equipaggio più diversificato inizia a sognare scene di mondi colonizzati e poi abbandonati dagli Iconiani.
\item[Il Giuramento di Protezione] L'equipaggio riceve visioni della caduta dell'Impero Iconiano, specificamente scene di evacuazioni.
\item[La Danza delle Stelle] I sistemi di navigazione mostrano stelle che si muovono in pattern coordinati, come un antico codice iconiano.
\end{description}
\end{tcolorbox}

\section{Personaggi Chiave}

\begin{tcolorbox}[colback=militarylightgray!10, colframe=militaryred]
\begin{description}[leftmargin=*]
\item[Comandante Aleksander Novak] Umano caldosiano, veterano pragmatico con esperienza diplomatica
\item[Tenente Comandante Sadira Veshi] Prima ufficiale anfibio-umanoide con capacità telepatiche limitate
\item[Maggiore Chen] Analista tattico esperto con conoscenza approfondita delle tecnologie antiche
\item[Maggiore Kron] Capo sicurezza tellarite, combattivo ma leale alla missione
\item[Tenente Sovar] Ufficiale scientifico romulano rifugiato, sotto costante sospetto
\end{description}
\end{tcolorbox}

\section{Principi Narrativi per Risoluzioni Diceless}

\begin{tcolorbox}[colback=militarylightgray!10, colframe=militaryred]
\begin{itemize}[leftmargin=*]
\item \textbf{Competenza}: Se un personaggio ha la competenza appropriata, riesce nell'azione base
\item \textbf{Complessità}: Per azioni complesse, richiedere descrizioni dettagliate dell'approccio
\item \textbf{Conseguenze}: Anche i successi hanno costi o complicazioni
\item \textbf{Collaborazione}: Due personaggi che lavorano insieme possono superare sfide più difficili
\item \textbf{Conflitto}: Le scelte contrastanti generano tensione narrativa senza necessità di dadi
\end{itemize}
\end{tcolorbox}

\section{Capacità Uniche della Nave}

\begin{tcolorbox}[colback=militarylightgray!10, colframe=militaryred]
\begin{description}[leftmargin=*]
\item[Intercettazione Teletrasporto] Può reindirizzare teletrasporti nemici (3 QDP)
\item[Blocco di Sicurezza] Può impedire tentativi di teletrasporto spendendo QDP
\item[Sistemi Difensivi Adattivi] Tecnologie derivate da diverse culture, adattabili a minacce sconosciute
\item[Conoscenza della Shackleton Expanse] Vantaggio negoziale nelle trattative con altre navi
\end{description}
\end{tcolorbox}

\section{Dilemmi Tattici}

\begin{tcolorbox}[colback=militarylightgray!10, colframe=militaryred]
\begin{itemize}[leftmargin=*]
\item Neutralizzare completamente l'artefatto o tentare di sfruttarne il potenziale difensivo
\item Condividere informazioni con altre navi o mantenere vantaggio informativo
\item Prioritizzare la protezione di alcune colonie rispetto ad altre
\item Fidarsi delle intuizioni del personale non-umano (Veshi telepatica, romulano rifugiato)
\item Rispettare le direttive della Coalizione o agire autonomamente per il bene maggiore
\end{itemize}
\end{tcolorbox}

\section{Meccanica del Teletrasporto in Presenza di Interferenze Iconiane}

\begin{tcolorbox}[colback=militarylightgray!10, colframe=militaryred]
\begin{itemize}[leftmargin=*]
\item \textbf{Costo Base}: 3 QDP
\item \textbf{Capacità Unica}: "Intercettazione Teletrasporto" (3 QDP)
\item \textbf{Difesa}: Può tentare di bloccare teletrasporti ostili spendendo QDP
\item \textbf{Bonus}: +1 se si utilizzano QDP Beta per precisione tattica
\item \textbf{Risoluzione}: 2d6 + modificatori, dove 10+ è successo completo
\end{itemize}
\end{tcolorbox}

\section{Opzioni per la Fase Finale}

\begin{tcolorbox}[colback=militarylightgray!10, colframe=militaryred]
\begin{description}[leftmargin=*]
\item[Neutralizzazione Completa] (6 QDP Beta + 2 QDP Alpha) Disabilitare permanentemente l'artefatto
\item[Scudo Difensivo] (4 QDP Beta + 4 QDP Alpha) Bloccare selettivamente i portali verso i mondi della Coalizione
\item[Riconfigurazione Tattica] (5 QDP di ogni tipo) Riprogrammare l'artefatto come asset difensivo
\end{description}
\end{tcolorbox}

\section{Gestione delle Scene Narrative}

\begin{tcolorbox}[colback=militarylightgray!10, colframe=militaryred]
\begin{itemize}[leftmargin=*]
\item Strutturare le scene con chiare opzioni e conseguenze
\item Per decisioni ad alta tensione, offrire ai giocatori opzioni con risultati prevedibili
\item Utilizzare il principio di "sì, ma..." permettendo successi con complicazioni
\item Per conflitti di priorità tra personaggi, incoraggiare il dibattito in-character
\item Stabilire countdown visibili per situazioni urgenti (es. attivazione dell'artefatto)
\end{itemize}
\end{tcolorbox}

\section{Equipaggio Multiculturale}

\begin{tcolorbox}[colback=militarylightgray!10, colframe=militaryred]
\begin{itemize}[leftmargin=*]
\item \textbf{Umani Caldosiani}: Adattati a mondi freddi, pragmatici e conservatori nelle risorse
\item \textbf{Veshi}: Anfibio-umanoidi con capacità telepatiche limitate e tradizione di navigazione
\item \textbf{Exofughi Tellariti}: Comunità separatiste con forte etica del dibattito costruttivo
\item \textbf{Rifugiati Romulani}: Sopravvissuti alla catastrofe di Hobus con competenze tecnologiche
\item \textbf{Catena di comando flessibile}: Basata sull'expertise più che sul grado militare
\end{itemize}
\end{tcolorbox}

\section{Comandi del Bot}

\begin{tcolorbox}[colback=militarylightgray!10, colframe=militaryred]
\begin{description}[leftmargin=*]
\item[/register] Registrazione come Master
\item[/status] Stato corrente dell'avventura
\item[/scan] Eseguire una scansione dell'artefatto
\item[/qdp\_status] Vedere lo stato dei tuoi QDP
\item[/add\_qdp <tipo> <quantità> <descrizione>] Aggiungere QDP (alpha, beta, gamma)
\item[/spend\_qdp <tipo> <quantità> <descrizione>] Spendere QDP
\item[/resolve\_enigma\_piece <tipo>] Risolvere parte di un enigma
\item[/teleport <tavolo> <descrizione>] Tentare un teletrasporto
\item[/convergence\_roll <descrizione>] Eseguire un tiro di Convergenza nella fase finale
\item[/send\_message <tavolo> <messaggio>] Inviare un messaggio ad un altro tavolo
\end{description}
\end{tcolorbox}

\end{document}